\exercisesection{Holomorphic Functional Calculus}

In the exercises below, all algebras considered are over C, unless explicitly stated otherwise.

\begin{enumerate}
\def\labelenumi{\arabic{enumi})}
\tightlist
\item
  Let \(n \geqslant 1\) be an integer and \(A = \mathcal{L}(\mathbf{C}^{n})\) . Let \(\alpha \in C\) and \(x \in A\) an element whose matrix relative to the canonical basis of \(C^{n}\) is a Jordan matrix of order n and eigenvalue \(\alpha\) The spectrum of x is reduced to \(\alpha\) .
\end{enumerate}

\begin{enumerate}
\def\labelenumi{\alph{enumi})}
\tightlist
\item
  For \(f \in \mathcal{O}(\{\alpha\})\) compute the matrix relative to the canonical basis of \(\mathbf{C}^n\) which represents \(f(x)\) .
\end{enumerate}

Let \(m \geqslant 1\) be an integer. Let \(B_{m}\) the unital Banach algebra of jets of functions of class \(C^{m}\) in a neighborhood of \(\alpha\) endowed with the topology of \(C^{m}\) -uniform convergence (VAR, R2, §12).

\begin{enumerate}
\def\labelenumi{\alph{enumi})}
\setcounter{enumi}{1}
\item
  There exists a continuous morphism of unital algebras from \(B_{m-1}\) to A such that \(\varphi(z)=x\) , where \(z\in B_{m-1}\) is the jet of order m-1 at \(\alpha\) of the identity function of C.
\item
  There does not exist a continuous morphism \(\varphi\) of unital algebras from \(B_m\) to A such that \(\varphi(z) = x\) , where \(z \in B_m\) is the jet of the identity function of \(\mathbf{C}\) .
\end{enumerate}

\begin{enumerate}
\def\labelenumi{\arabic{enumi})}
\setcounter{enumi}{1}
\tightlist
\item
  Let A be a unital Banach algebra and G the group of invertible elements of A.
\end{enumerate}

\begin{enumerate}
\def\labelenumi{\alph{enumi})}
\item
  For an element \(x \in A\) to be of the form \(\exp(y)\) , where \(y \in A\) , it is necessary and sufficient that it be contained in a connected commutative subgroup of G.
\item
  For an element \(x \in \mathbf{A}\) to be of the form \(\exp(y)\) , where \(y \in \mathbf{A}\) , it suffices that 0 belong to the unbounded connected component of \(\mathbf{C} = \mathrm{Sp}_{\mathrm{A}}(x)\) .
\end{enumerate}

\begin{enumerate}
\def\labelenumi{\arabic{enumi})}
\setcounter{enumi}{2}
\tightlist
\item
  Let A be a unital Banach algebra, G the group of invertible elements of A, and \(G_{1}\) the identity component of G.
\end{enumerate}

\begin{enumerate}
\def\labelenumi{\alph{enumi})}
\item
  Let \(x \in \mathbf{A}\) and \(n\) an integer \(\geqslant 0\) such that \(x^n = e\) . Show that \(x \in \mathbf{G}_1\) . (The spectrum of \(x\) is finite. The set of \(\lambda \in \mathbf{C}\) such that \(y(\lambda) = \lambda x + e - \lambda e\) such that it is invertible has finite complement, hence is connected. Now, \(y(0) = e\) , \(y(1) = x\) .)
\item
  Assume A is commutative. Every element of G/G₁ distinct from the identity element has infinite order. (If \(x \in G\) and \(x^n \in G_1\) , one has \(x^n = \exp(y)\) with a \(y \in A\) by no. 10 of I, p.~78, whence ( \(x \exp(-y/n)^n = e\) . Apply a).)
\end{enumerate}

\begin{enumerate}
\def\labelenumi{\arabic{enumi})}
\setcounter{enumi}{3}
\item
  There exists a commutative Banach algebra, non-unital and non-zero, such that \(X(A)\) is compact.
\item
  Let A be a unital Banach algebra and \(a \in A\) . Let G be the group of invertible elements of A and \(G_1\) the connected component of the identity in G. The exponential spectrum of A is the set of \(\lambda \in \mathbf{C}\) such that \(x - \lambda e\) does not belong to \(G_1\) (cf.~prop. 12 of I, p.~79).
\end{enumerate}

\begin{enumerate}
\def\labelenumi{\alph{enumi})}
\item
  The exponential spectrum of a is closed; it contains the spectrum of a.
\item
  The boundary of the exponential spectrum of a is contained in the singular spectrum of a (exercise 16 of I, p.~158).
\item
  The exponential spectrum of a is the union of the spectrum of a and of certain bounded connected components of \(\mathbf{C}=\mathrm{Sp}(a)\) .
\end{enumerate}

\begin{enumerate}
\def\labelenumi{\alph{enumi})}
\setcounter{enumi}{3}
\tightlist
\item
  Let \(f\in \mathbf{C}[\mathrm{X}]\) be a polynomial. The image under \(f\) of the exponential spectrum of \(a\) contains the exponential spectrum of \(f(a)\) . In general, the inclusion is strict (use Fredholm theory, cf.~III, to appear).
\end{enumerate}

\begin{enumerate}
\def\labelenumi{\arabic{enumi})}
\setcounter{enumi}{5}
\tightlist
\item
  Let A be a commutative unital Banach algebra.
\end{enumerate}

\begin{enumerate}
\def\labelenumi{\alph{enumi})}
\item
  If j is an idempotent of A, one has \(\exp(2i\pi j)=1\) and \(\exp(i\pi j)=1-2j\) .
\item
  Let \(p \in A\) such that \(\exp(p) = 1\) . Then \(\mathrm{Sp}_{\mathrm{A}}(p)\) consists of a finite number of points of the form \(2in\pi\) , where \(n \in Z\) .
\item
  For \(\lambda \notin \mathrm{Sp}_{\mathrm{A}}(p)\) , we have:
\end{enumerate}

\[
\int_ {0} ^ {1} \exp (t (p - \lambda 1)) d t = (1 - \exp (- \lambda)) (\lambda 1 - p) ^ {- 1}.
\]

\begin{enumerate}
\def\labelenumi{\alph{enumi})}
\setcounter{enumi}{3}
\tightlist
\item
  Deduce that every point of \(\mathrm{Sp}_{\mathbf{A}}(p)\) is a simple pole of the resolvent \(\mathrm{R}(p,\lambda)\) , then that there exists an integer \(m \geqslant 0\) , integers \(n_{k} \in Z\) and pairwise orthogonal idempotents \(j_{k}\) , for \(1 \leqslant k \leqslant m\) , such that
\end{enumerate}

\[
p = 2 i \pi \sum_ {k} n _ {k} j _ {k}.
\]

\begin{enumerate}
\def\labelenumi{¶ \arabic{enumi})}
\setcounter{enumi}{6}
\tightlist
\item
  a) Let B be a complex Banach space and j a point of B such that \(\|j\| = 1\) . For every \(x \in B\) , the limit
\end{enumerate}

\[
\lim_{\substack{\alpha \to 0\\ \alpha >0}}\frac{1}{\alpha} (\| j + \alpha x\| -1)
\]

exists. Let \(\varphi(x)\) be this limit.

\begin{enumerate}
\def\labelenumi{\alph{enumi})}
\setcounter{enumi}{1}
\tightlist
\item
  Let
\end{enumerate}

\[
\psi (x) = \sup_{\substack{\varepsilon \in \mathbf{C}\\ |\varepsilon | = 1}}\varphi (\varepsilon x).
\]

The function \(\psi\) is a semi-norm bounded above by the norm of B. One has

\[
\psi (x) = \sup _ {f} | f (x) |,
\]

where \(f\) ranges over the set of elements of the dual of B such that \(\|f\| = f(j) = 1\) .

\begin{enumerate}
\def\labelenumi{\alph{enumi})}
\setcounter{enumi}{2}
\tightlist
\item
  One now assumes that B is a Banach algebra and that \(j = 1\) is a unit element of B. One has:
\end{enumerate}

\[
\varphi (x) = \lim_{\substack{\alpha \to 0\\ \alpha >0}}\frac{1}{\alpha}\log (\| \exp (\alpha x)\|) = \sup_{\alpha >0}\frac{1}{\alpha}\log (\| \exp (\alpha x)\|).
\]

(Observe that \(\log(\|\exp((\alpha+\alpha')x)\|) \leqslant \log(\|\exp(\alpha x)\|) + \log(\|\exp(\alpha' x)\|).\) )

\begin{enumerate}
\def\labelenumi{\alph{enumi})}
\setcounter{enumi}{3}
\tightlist
\item
  Using c), show that \(\|\exp(\lambda x)\| \leqslant \exp(|\lambda|\psi(x))\) . By integrating \(\lambda \mapsto \lambda^{-2} \exp(\lambda x)\) over the circle determined by \(|\lambda| = \varrho\) , deduce that
\end{enumerate}

\[
\| x \| \leqslant \varrho^ {- 1} \exp (\varrho \psi (x)),
\]

whence, setting \(\varrho = \psi(x)^{-1}\) , that \(\|x\| \leqslant e\psi(x)\) , where \(e = \exp(1)\) .

\begin{enumerate}
\def\labelenumi{\alph{enumi})}
\setcounter{enumi}{4}
\tightlist
\item
  For every \(x\) nonzero element of B, there exists a linear form \(f \in \mathrm{B}'\) such that \(f(1) = \| f \| = 1\) and \(f(x) \neq 0\) .
\end{enumerate}

\begin{enumerate}
\def\labelenumi{\arabic{enumi})}
\setcounter{enumi}{7}
\item
  Let A be a unital Banach algebra, \(x \in A\) , and B the subalgebra of A generated by x. For \(\mathrm{Sp}_{\mathrm{A}}(x)\) to be a finite set all of whose elements are poles of the resolvent of x, it is necessary and sufficient that B be of finite dimension. (To show that the condition is necessary, if \(\mathrm{R}(x, \lambda) = \sum_{i=1}^{p} \mathrm{F}_{i}(\lambda) a_{i}\) , where \(a_{i} \in A\) and where \(F_{i}\) is a scalar rational function, the expansion of \(\mathrm{R}(x, \lambda)\) around \(\lambda = \infty\) shows that the \(x^{n}\) are linear combinations of the \(a_{i}\) .)
\item
  Let A be a unital Banach algebra, \(x \in A\) , U an open neighborhood of \(\mathrm{Sp}_{\mathrm{A}}(x)\) having only a finite number of connected components \(U_{i}\) ( \(i \in I\) ), and \(f \in \mathcal{O}(\mathrm{U})\) . For \(f(x) = 0\) , it is necessary and sufficient that the following conditions be satisfied:
\end{enumerate}

\begin{enumerate}
\def\labelenumi{(\roman{enumi})}
\item
  for every i such that \(\mathrm{Sp}_{\mathrm{A}}(x)\cap\mathrm{U}_{i}\) is infinite or contains a point which is not a pole of the resolvent, one has \(f|U_{i}=0\) ;
\item
  every pole of order \(p\) of the resolvent is a zero of order \(\geqslant p\) of \(f\) (For necessity of (ii), argue as in Exerc. 8.)
\end{enumerate}

\begin{enumerate}
\def\labelenumi{\arabic{enumi})}
\setcounter{enumi}{9}
\tightlist
\item
  Let \(\Delta\) the disk \(|z| \leqslant 1\) in C. Let A be the space of functions f that are continuous in \(\Delta\) and holomorphic in the interior of \(\Delta\) and such that there exists a sequence \((c_{n})\) summable of complex numbers satisfying
\end{enumerate}

\[
f (z) = \sum_ {n \geqslant 0} c _ {n} z ^ {n}
\]

for all \(z \in \Delta\) . We then denote by \(\|f\| = \sum_{n}|c_{n}|\) .

\begin{enumerate}
\def\labelenumi{\alph{enumi})}
\item
  Equipped with the multiplication \((f,g)\mapsto fg\) , the space A is a unital Banach algebra and z is a topological generator of A.
\item
  The space \(\mathsf{X}(\mathsf{A})\) is identified with \(\Delta\) ; if \(f \in A\) and if g is holomorphic in an open neighborhood of \(f(\Delta)\) , then \(g \circ f \in A\) .
\end{enumerate}

\begin{enumerate}
\def\labelenumi{\arabic{enumi})}
\setcounter{enumi}{10}
\item
  Let A be a commutative unital Banach algebra and \(x \in A\) . Let S be a subset of X(A) closed for the Jacobson topology, and let f be a complex holomorphic function in a neighborhood of \((\mathcal{G}x)(\mathrm{S})\) . There exists \(y \in A\) such that \(Gy = f \circ Gx\) on S. (Let I be the intersection of the kernels of the \(\chi \in S\) . Use the functional calculus in A/I.)
\item
  Let U be a nonempty open set in C. Show that there exists no norm on \(\mathcal{O}(\mathrm{U})\) for which \(\mathcal{O}(\mathrm{U})\) is a Banach algebra. (Using th. 1 of I, p.~29, show that the functions of \(\mathcal{O}(\mathrm{U})\) would be bounded.)
\item
  Let \(n\in \mathbf{N}\)
\end{enumerate}

\begin{enumerate}
\def\labelenumi{\alph{enumi})}
\tightlist
\item
  A subset V of \(C^{n}\) is polynomially convex if and only if, for every \(x \in C^{n} - V\) , there exists an entire function \(f: C^{n} \to C\) such that
\end{enumerate}

\[
| f (x) | > \sup _ {y \in V} | f (y) |.
\]

\begin{enumerate}
\def\labelenumi{\alph{enumi})}
\setcounter{enumi}{1}
\item
  Let \(K \subset R^{n}\) a compact set. Then K is polynomially convex in \(C^{n}\) .
\item
  Let V be a polynomially convex compact subset of \(C^{n}\) and \(f \in \mathcal{O}(V)\) . Then the graph of f is polynomially convex in \(C^{n+1}\) . (Use the Oka--Weil theorem.)
\end{enumerate}

\begin{enumerate}
\def\labelenumi{\arabic{enumi})}
\setcounter{enumi}{13}
\item
  Let K be a compact subset of C and let P (resp. Q) be the closure in \(\mathcal{O}(\mathrm{K})\) of the set of germs of polynomial (resp. rational) functions holomorphic in a neighborhood of K. Show that \(\mathrm{P} = \mathrm{Q}\) if and only if the set \(\mathbf{C} - \mathbf{K}\) is connected.
\item
  Let U be the set of complex numbers of modulus 1. Let \(K_{1}\) , \(K_{2}\) be the compact subsets of \(C^{2}\) given by
\end{enumerate}

\[
\mathrm{K} _ {1} = \mathbf {U} \times \{0 \}, \quad \mathrm{K} _ {2} = \{(z, \overline {{z}}) \mid z \in \mathbf {U} \}.
\]

\begin{enumerate}
\def\labelenumi{\alph{enumi})}
\item
  Show that the set of germs in a neighborhood of \(K_{1}\) of polynomial functions on C with complex coefficients is not dense in \(\mathcal{O}(K_{1})\) .
\item
  Show that the set of germs in a neighborhood of \(K_{2}\) of polynomial functions on C with complex coefficients is dense in \(\mathcal{O}(K_{2})\) .
\end{enumerate}

\begin{enumerate}
\def\labelenumi{\arabic{enumi})}
\setcounter{enumi}{15}
\tightlist
\item
  Let A be a unital Banach algebra, \(x \in A\) , and k an integer \(\geqslant 0\) . Let U be the set of complex numbers of modulus 1.
\end{enumerate}

\begin{enumerate}
\def\labelenumi{\alph{enumi})}
\tightlist
\item
  Suppose that \(\mathrm{Sp}_{\mathrm{A}}(x)\subset\mathbf{U}\) and that \(\|\mathrm{R}(x,\lambda)\|=\mathrm{O}((1-|\lambda|)^{k})\) when \(|\lambda|\) tends to 1. Show that \(\|x^{n}\|=\mathrm{O}(|n|^{k})\) when n tends to \(\pm\infty\) . (For every \(\delta>1\) , show that
\end{enumerate}

\[
2 \pi \| x ^ {n} \| \leqslant \mathrm{M} \frac {\delta^ {n}}{(\delta - 1) ^ {k}} + \mathrm{M} \frac {\delta^ {- n}}{(1 - \delta^ {- 1}) ^ {k}}
\]

where M is independent of n and \(\delta\) . Take \(\delta = n/(n - k)\) when n tends to \(+\infty\) , \(\delta = n/(n + k)\) when n tends to \(-\infty\) .

\begin{enumerate}
\def\labelenumi{\alph{enumi})}
\setcounter{enumi}{1}
\tightlist
\item
  Suppose that \(\|x^{n}\|=\mathrm{O}(|n|^{k})\) when n tends to \(\pm\infty\) . Show that \(\operatorname{Sp}_{\mathbf{A}}(x)\subset\mathbf{U}\) and that \(\|\operatorname{R}(\lambda,x)\|=\operatorname{O}((1-|\lambda|)^{k+1})\) when \(|\lambda|\) tends to 1. (We have \(\varrho(x)=\varrho(x^{-1})=1\) , hence \(\operatorname{Sp}_{\mathbf{A}}(x)\subset\mathbf{U}\) . For \(|\lambda|<1\) ,
\end{enumerate}

\[
\mathrm{R} (\lambda , x) = - x ^ {- 1} \left(1 + \lambda x ^ {- 1} + \lambda^ {2} x ^ {- 2} + \lambda^ {3} x ^ {- 3} + \dots\right).
\]

Deduce an upper bound for \(\|\mathrm{R}(\lambda,x)\|\) , and argue similarly if \(|\lambda|>1\) .

\begin{enumerate}
\def\labelenumi{\arabic{enumi})}
\setcounter{enumi}{16}
\tightlist
\item
  Let \(X_{1}\) (resp. \(X_{2}\) ) the set of \((z_{1}, z_{2}) \in \mathbf{C}^{2}\) such that
\end{enumerate}

\[
\frac {1}{2} \leqslant | z _ {1} | ^ {2} + | z _ {2} | ^ {2} \leqslant 1, \quad \text { resp. } \quad | z _ {1} | ^ {2} + | z _ {2} | ^ {2} \leqslant 1.
\]

Let \(X = X_{1} \cup \{(0,0)\} \subset \mathbf{C}^{2}\) . Let \(A = \mathcal{C}(X)\) , and \(a_{1}, a_{2}\) the restrictions to X of the coordinate functions on \(C^{2}\) .

\begin{enumerate}
\def\labelenumi{\alph{enumi})}
\item
  The simultaneous spectrum \(\mathrm{Sp}_{\mathrm{A}}^{2}(a_{1},a_{2})\) is equal to X.
\item
  Construct two distinct continuous unital morphisms from \(\mathcal{O}(\mathrm{X})\) to A sending \(z_{1}\) to \(a_{1}\) and \(z_{2}\) to \(a_{2}\) . (Use the following theorem: for every function \(f\) holomorphic in a neighborhood of \(\mathrm{X}_{1}\) , there exists a function g holomorphic in a neighborhood of \(X_{2}\) which coincides with f in a neighborhood of \(X_{1}\) (cf. I, p. 168, exerc. 14).
\end{enumerate}

\begin{enumerate}
\def\labelenumi{¶ \arabic{enumi})}
\setcounter{enumi}{17}
\tightlist
\item
  Let A be a commutative unital Banach algebra.
\end{enumerate}

\begin{enumerate}
\def\labelenumi{\alph{enumi})}
\tightlist
\item
  For any finite family I of elements of A, let S(I) ⊂ C \(^{I}\) the simultaneous spectrum of I. If I′ extends I, the canonical surjection pr \(_{I,I'}\) of C \(^{I'}\) on C \(^{I}\) is such that pr \(_{I,I'}\) (S(I')) = S(I), whence an injective morphism \(\varphi_{II'}\) of \(\mathcal{O}(S(I))\) into \(\mathcal{O}(S(I'))\) . Let \(\mathcal{O}(X(A))\) the inductive limit of the \(\mathcal{O}(S(I))\) with respect to the \(\varphi_{II'}\) (we index the finite families of elements of A by the finite subsets of A × N). The canonical maps
\end{enumerate}

\[
\varphi_ {\mathrm{I}} \colon \mathcal {O} (\mathrm{S} (\mathrm{I})) \longrightarrow \mathcal {O} (\mathrm{X} (\mathrm{A}))
\]

are injective.

We shall equip \(\mathcal{O}(\mathrm{X(A)})\) of the inductive limit topology of those of \(\mathcal{O}(\mathrm{S(I)})\) . For every \(a \in \mathrm{A}\) the coordinate function in a neighborhood of \(\mathrm{S}(a) \subset \mathbf{C}\) defines an element \(z_a\) of \(\mathcal{O}(\mathrm{X(A)})\) .

\begin{enumerate}
\def\labelenumi{\alph{enumi})}
\setcounter{enumi}{1}
\tightlist
\item
  The continuous surjection \(\mathsf{X}(\mathsf{A}) \to \mathsf{S}(\mathsf{I})\) defines a continuous morphism \(\mathcal{C}(\mathsf{S}(\mathsf{I})) \to \mathcal{C}(\mathsf{X}(\mathsf{A}))\) . The morphisms \(\mathcal{O}(\mathsf{S}(\mathsf{I})) \to \mathcal{C}(\mathsf{S}(\mathsf{I})) \to \mathcal{C}(\mathsf{X}(\mathsf{A}))\) define a continuous morphism:
\end{enumerate}

\[
\varphi \colon \mathcal {O} (\mathrm{X} (\mathrm{A})) \longrightarrow \mathcal {C} (\mathrm{X} (\mathrm{A})).
\]

This morphism is not injective in general.

\begin{enumerate}
\def\labelenumi{\alph{enumi})}
\setcounter{enumi}{2}
\item
  There exists a unique continuous unital morphism \(\psi\) of \(\mathcal{O}(\mathsf{X}(\mathsf{A}))\) to A such that \(\psi(z_a) = a\) for all \(a \in \mathsf{A}\) (Use th. 1 of I, p.~51.) Its composite with the Gelfand transformation is \(\varphi\) .
\item
  Let E be a separated locally convex space, U a weakly open subset of E, and \(f: U \to C\) a function. We say that f is weakly holomorphic if, for every \(u \in U\) there exists a weakly open neighborhood \(V_u\) of u, a closed vector subspace \(F_u\) of finite codimension in E, and a holomorphic function \(g_u\) on \(p_u(V_u)\) , where \(p_u\) denotes the canonical mapping of E onto \(E/F_u\) , such that \(f|V_u = g_u \circ p_u\) .
\end{enumerate}

Show that if \(f\) is weakly holomorphic, and if K is a weakly compact subset of U, there exists a weakly open neighborhood \(V_K\) of K in U, a closed vector subspace \(F_K\) of finite codimension in E, and a holomorphic function \(g_K\) on \(p_K(V_K)\) such that \(f|V_K = g_K \circ p_K\) , where \(p_K\) is the canonical mapping of E onto \(E/F_K\) .

\begin{enumerate}
\def\labelenumi{\alph{enumi})}
\setcounter{enumi}{4}
\tightlist
\item
  Show that \(\mathcal{O}(\mathsf{X}(\mathsf{A}))\) is the inductive limit of the \(\mathcal{O}(\mathsf{S})\) , where S ranges over the set of weakly open neighborhoods of \(\mathsf{X}(\mathsf{A})\) in the dual of A, and where \(\mathcal{O}(\mathsf{S})\) denotes the algebra of weakly holomorphic functions in S, endowed with the topology of compact convergence.
\end{enumerate}

